\begin{figure}[t]
\centering
\begin{tikzpicture}[x=0.85cm,y=0.72cm,>=Stealth,font=\small]
  \draw[->] (0,1) -- (8.2,1) node[right] {$t$};
  \foreach \y/\i/\start/\finish in {0/1/1/4.2,-1.4/2/2.5/6.5,-2.8/3/0.6/5.2} {
    \node[left] at (0,\y) {$x_{\i}$};
    \draw[gray!60] (0,\y)--(\start,\y);
    \draw[very thick,blue!65!black] (\start,\y)--(\finish,\y);
    \draw[dashed,thick] (\finish,\y)--(7.6,\y);
    \fill (\start,\y) circle(1.5pt);
    \fill (\finish,\y) circle(1.5pt);
    \node[above] at (\start,\y) {$\sigma_{\i}$};
    \node[above] at (\finish,\y) {$T_{\i}$};
  }
  \node[align=center] at (2.2,-4) {solid: active factor};
  \node[align=center] at (6,-4) {dashed: frozen factor};
\end{tikzpicture}
\caption{The multipoint event uses each point's own activation and target
stop. The solid intervals may overlap in any order. After $T_i$ the
factor for $x_i$ is constant, and no further survival of that point is
required. The axis is schematic physical time, not the individual
conformal-radius clocks.}
\label{fig:individual-stops}
\end{figure}
